\section{1.06 building first binary}
\label{sec:1.06 building first binary}

% Section 01.06: Building Your First Embedded Binary

Now that we have our toolchain set up and understand the project structure, let's build and analyze a real embedded binary. This section will guide you through the complete build process with detailed explanations of each step.

\subsection{Preparing a Complete Example}
\label{sec:complete-example}

Let's create a minimal but complete embedded project that we can actually build and analyze:
\begin{minted}{bash}
$ cargo new --lib efr32mg24-minimal
$ cd efr32mg24-minimal
$ mkdir examples
$ cat > examples/minimal.rs << 'EOF'
#![no_std]
#![no_main]

use core::panic::PanicInfo;

#[panic_handler]
fn panic(_info: &PanicInfo) -> ! {
    loop {}
}

#[no_mangle]
pub extern "C" fn main() -> ! {
    loop {
        // Main application loop
        cortex_m::asm::nop();
    }
}
EOF
\end{minted}

Add dependencies to \filepath{Cargo.toml}:\

\begin{minted}{toml}
[package]
name = "efr32mg24-minimal"
version = "0.1.0"
edition = "2021"

[profile.dev]
opt-level = "z"
debug = true
overflow-checks = false  # Critical for embedded

[profile.release]
opt-level = "z"
debug = true
overflow-checks = false
panic = "abort"
lto = true

codegen-units = 1

[dependencies]
cortex-m = "0.7"
cortex-m-rt = "0.7"

# Development dependencies
[dev-dependencies]
cortex-m-semihosting = "0.5"
\end{minted}

\subsection{Reference: Complete Project Example}
\label{subsec:reference-complete-project}

For a comprehensive step-by-step guide with all project files (build.rs, Cargo.toml, memory.x, etc.), see Section \ref{sec:complete-project-example}.

This section focuses on understanding the build process and analyzing the resulting binary.

\subsection{Understanding the Build Process}
\label{sec:understanding-build}

When you run \codeinline{cargo build --release}, here's what happens:
\begin{remark}
    This diagram shows the build pipeline for embedded Rust:
\end{remark}

\begin{figure}[H]
    \centering
    \begin{tikzpicture}[
        node/.style={rectangle, draw=black, thick, minimum width=2.5cm, minimum height=1cm},
        arrow/.style={-Latex, thick}
    ]

        % Nodes
        \node[node, fill=rustblue!20] (source) at (0, 0) {Source Code};
        \node[node, fill=yellow!20] (rustc) at (3, 0) {rustc};
        \node[node, fill=orange!20] (llvm) at (6, 0) {LLVM};
        \node[node, fill=green!20] (obj) at (9, 0) {Object File};
        \node[node, fill=red!20] (linker) at (9, -2) {Linker};
        \node[node, fill=purple!20] (bin) at (6, -2) {Binary};

        % External files
        \node[font=\small] (linker-script) at (9, 0.8) {linker.ld};
        \node[font=\small] (startup) at (9, 1.2) {startup_s};

        % Arrows
        \draw[arrow] (source) -- (rustc);
        \draw[arrow] (rustc) -- (llvm);
        \draw[arrow] (llvm) -- (obj);
        \draw[arrow] (obj) -- (linker);
        \draw[arrow] (linker) -- (bin);
        \draw[arrow, dashed] (linker-script) -- (linker);
        \draw[arrow, dashed] (startup) -- (linker);

        % Labels
        \node[font=\small, below=0.2cm of source] {Human};
        \node[font=\small, below=0.2cm of rustc] {Targeting};
        \node[font=\small, below=0.2cm of llvm] {ARM};
        \node[font=\small, below=0.2cm of obj] {Cortex-M33};
        \node[font=\small, below=0.2cm of bin] {thumbv8m};
    \end{tikzpicture}
    \caption{Rust embedded build pipeline}
    \label{fig:build-pipeline}
\end{figure}

\subsection{Compile Step}
\begin{minted}{bash}
$ cargo build --release --example minimal
   Compiling autocfg v1.1.0
   Compiling proc-macro2 v1.0.56
   Compiling cortex-m v0.7.7
   Compiling quote v1.0.27
   Compiling cortex-m-rt v0.7.3
   Compiling efr32mg24-minimal v0.1.0
    Finished release [optimized] target(s) in 1.24s
\end{minted}

Notice that cargo automatically downloads and compiles dependencies. For embedded,
this includes startup code and linker scripts via \codeinline{cortex-m-rt}.

\subsection{Binary Analysis}
\label{sec:binary-analysis}

Once built, let's examine the binary:

\subsection{Size Analysis}
\begin{minted}{bash}
$ cargo size --release --example minimal
   text    data     bss     dec     hex filename
     80       0       0      80      50 minimal
\end{minted}

What do these sections mean?
\begin{itemize}
    \item \textbf{text}: The actual machine code (read-only)
    \item \textbf{data}: Initialized static variables
    \item \textbf{bss}: Uninitialized static variables (zeroed at startup)
    \item \textbf{dec}: Total bytes (decimal)
    \item \textbf{hex}: Total bytes (hexadecimal)
\end{itemize}

Our minimal program is 80 bytes total. Let's see what's inside:

\subsection{Disassembly}
\begin{minted}{bash}
$ cargo objdump --release --example minimal -- -d
\end{minted}

Output (simplified):
\begin{minted}{text}
00000000 <main>:
       0: b580            push    {r7, lr}
       2: f000 f802       bl      0xa <_ZN4core4...
       6: bf00            nop
       8: e7fe            b.n     0x8 <main+0x8>

0000000a <HardFault>:
       a: e7fe            b.n     0xa <HardFault>
\end{minted}

Interpretation:
\begin{itemize}
    \item \codeinline{push {r7, lr}}: Save registers
    \item \codeinline{bl}: Branch with link (call function)
    \item \codeinline{nop}: No operation (our loop body)
    \item \codeinline{b.n}: Branch to self (infinite loop)
    \item \codeinline{HardFault}: Default exception handler
\end{itemize}

\subsection{Listing Sections}
\begin{minted}{bash}
$ cargo objdump --release --example minimal -- -h
\end{minted}

Shows memory layout:
\begin{minted}{text}
Idx Name          Size      VMA       LMA       File off  Algn
  0 .text         00000014  00000000  00000000  00000034  2**2
                  CONTENTS, ALLOC, LOAD, READONLY, CODE
  1 .data         00000000  20000000  20000000  00000034  2**0
                  CONTENTS, ALLOC, LOAD, DATA
  2 .bss          00000000  20000000  20000000  00000034  2**0
                  ALLOC
\end{minted}

\textbf{Explanation}:
\begin{itemize}
    \item \codeinline{.text} starts at 0x00000000 (flash memory)
    \item \codeinline{.data/.bss} start at 0x20000000 (SRAM)
    \item Flash: 1,536 KB on EFR32MG24
    \item SRAM: 256 KB on EFR32MG24
\end{itemize}

\subsection{Build Optimization Comparison}
\label{sec:optimization-comparsion}

Build with different optimization levels:

\begin{table}[H]
    \centering
    \begin{tabular}{l c c c c}
        \toprule
        \textbf{Profile} & \textbf{text} & \textbf{data} & \textbf{bss} & \textbf{Total} \\
        \midrule
        dev (opt-level='z') & 124 & 0 & 0 & 124 \\
        release (opt-level='z') & 80 & 0 & 0 & 80 \\
        release (opt-level='s') & 76 & 0 & 0 & 76 \\
        release (opt-level=3) & 68 & 0 & 0 & 68 \\
        \bottomrule
    \end{tabular}
    \caption{Size comparison by optimization level}
    \label{tab:optimization-size}
\end{table}

Optimization levels:
\begin{itemize}
    \item \codeinline{'z'}: Optimize for size aggressively
    \item \codeinline{'s'}: Optimize for size
    \item \codeinline{'3'}: Optimize for speed (may reduce size too)
\end{itemize}

Even the unoptimized debug build is only 124 bytes!

\subsection{Linking and Linker Scripts}
\label{sec:linking}

The linker script is crucial in embedded systems. It tells the linker where to place code and data in memory.
\subsection{Memory Layout}
Create \filepath{memory.x}:
\begin{minted}{text}
MEMORY
{
  /* Flash memory for code and constants */
  FLASH (rx) : ORIGIN = 0x00000000, LENGTH = 1536K

  /* RAM for variables and stack */
  RAM (rwx)  : ORIGIN = 0x20000000, LENGTH = 256K
}
\end{minted}

This matches the EFR32MG24B220 memory map:
\begin{itemize}
    \item Flash: 1,536 KB starting at 0x0000_0000
    \item SRAM: 256 KB starting at 0x2000_0000
\end{itemize}
\subsection{Examining the Linker Output}
\begin{minted}{bash}
$ cargo build --release --example minimal --verbose 2>&1 | grep "LINK"
\end{minted}

Shows the linker command:
\begin{minted}{text}
"rust-lld" "-flavor" "gnu"
  "-Tlink.x"
  "-Tmemory.x"
  "-o" ".../target/release/examples/minimal"
  ...
\end{minted}

\textbf{Key flags}:
\begin{itemize}
    \item \codeinline{rust-lld}: LLVM linker for embedded targets
    \item \codeinline{-Tlink.x}: Standard Cortex-M linker script
    \item \codeinline{-Tmemory.x}: Our custom memory layout
\end{itemize}

\subsection{Creating a Flashable Binary}
\label{sec:flashable-binary}

The ELF binary contains debug symbols and metadata. For flashing, we need a raw binary:
\subsection{Convert to Binary Format}
\begin{minted}{bash}
$ cargo objcopy --release --example minimal -- -O binary minimal.bin
\end{minted}

Check sizes:
\begin{minted}{bash}
$ ls -lh target/thumbv8m.main-none-eabihf/release/
-rwxr-xr-x  80K minimal          # ELF with symbols
-rwxr-xr-x  80B minimal.bin     # Raw binary for flash
\end{minted}

The binary is much smaller because it contains only the raw machine code without debug symbols, section headers, etc.
\subsection{Hex Format for Some Tools}
Some tools prefer Intel hex format:
\begin{minted}{bash}
$ cargo objcopy --release --example minimal -- -O ihex minimal.hex
\end{minted}

\subsection{Automated Build Script}
\label{sec:automated-build}

Create a rust file, \filepath{build.rs}, in your project root:
\begin{minted}{rust}
use std::env;
use std::fs;
use std::path::PathBuf;

fn main() {
    // Tell cargo to rerun build script if memory.x changes
    println!("cargo:rerun-if-changed=memory.x");

    // Copy memory.x to output directory
    let out = &PathBuf::from(env::var_os("OUT_DIR").unwrap());
    fs::copy("memory.x", out.join("memory.x")).unwrap();

    println!("cargo:rustc-link-search={}", out.display());
}
\end{minted}

This ensures the linker script is available during linking.

\subsection{Build Automation with Justfile}
\label{sec:justfile}

For complex projects, use \structure{just} (a command runner):
\begin{minted}{text}
# Justfile - Build automation

# Build release binary
build:
    cargo build --release --example minimal

# Show binary size
size:
    cargo size --release --example minimal

# Disassemble
disasm:
    cargo objdump --release --example minimal -- -d

# Show section info
sections:
    cargo objdump --release --example minimal -- -h

# Convert to binary
binary:
    cargo objcopy --release --example minimal -- -O binary minimal.bin

# Run all analysis
analyze: build size disasm sections binary
    @echo "Analysis complete!"
\end{minted}

Usage:
\begin{minted}{bash}
$ just build
$ just size
$ just analyze
\end{minted}
\subsection{Key Takeaways}
\label{sec:ch1-key-takeaways}

\begin{keytakeaways}
    \item Embedded Rust binaries are incredibly small (our minimal example is 80 bytes)
    \item Optimization settings significantly impact size (from 124 to 68 bytes)
    \item The build process involves: compile → link → optionally convert format
    \item Linker scripts control memory layout and are critical for embedded
    \item Cargo provides powerful tools for analyzing binaries (size, objdump, objcopy)
    \item The ELF format contains metadata; for flashing, use binary or hex format
    \item Scripting your build workflow improves productivity
\end{keytakeaways}

\begin{exercise}[Easy]
    Create a project template with a complete build script that:
    \begin{itemize}
        \item Builds release binary
        \item Shows size information
        \item Creates both .bin and .hex files
        \item Runs objdump analysis
    \end{itemize}
\end{exercise}

\begin{exercise}[Medium]
    Compare binaries built with different optimization levels:
    \begin{itemize}
        \item What differences do you see in objdump output?
        \item How much code can the linker remove from our minimal example?
        \item Where does libc start inserting startup code?
    \end{itemize}
\end{exercise}

\begin{exercise}[Hard]
    Write a custom linker script that:
    \begin{itemize}
        \item Places the vector table at a specific address
        \item Creates a separate section for time-critical code
        \item Reserves space at the end of RAM for a stack canary
        \item Adds checksums for firmware verification
    \end{itemize}
    Then verify it works with objdump.
\end{exercise}
